% Môn: Toán
% Lớp: 11
% Chương: Chương I. Hàm số lượng giác và phương trình lượng giác
% Bài: Bài 1. Giá trị lượng giác của góc lượng giác
% Loại: Dạng mẫu
% File riêng, không trộn vào de.tex
% Định nghĩa lệnh Lý thuyết — dùng khi biên dịch lt.tex bằng TeX.
% App web đọc lt.tex trực tiếp (bỏ \input file này).
% Trong lt.tex, từ thư mục bài:
%   % Định nghĩa lệnh Lý thuyết — dùng khi biên dịch lt.tex bằng TeX.
% App web đọc lt.tex trực tiếp (bỏ \input file này).
% Trong lt.tex, từ thư mục bài:
%   % Định nghĩa lệnh Lý thuyết — dùng khi biên dịch lt.tex bằng TeX.
% App web đọc lt.tex trực tiếp (bỏ \input file này).
% Trong lt.tex, từ thư mục bài:
%   \input{../../../../_lenh/lythuyet.tex}
% từ gốc repo:
%   \input{ngan-hang/_lenh/lythuyet.tex}
\makeatletter
\providecommand{\indam}[1]{\textbf{#1}}
\providecommand{\chuy}[1]{\par\smallskip\noindent\textit{#1}\par}
\providecommand{\luuy}[1]{\par\smallskip\noindent\textbf{Lưu ý.} #1\par}
\providecommand{\ghichu}[1]{\par\smallskip\noindent\textbf{Ghi chú.} #1\par}
\providecommand{\vidu}[1]{\par\smallskip\noindent\textbf{Ví dụ.} #1\par}
\providecommand{\Blythuyet}{\subsection*{Lý thuyết}}
% Hai cột: kiến thức | hình
\providecommand{\haicotchay}[2]{%
  \par\medskip\noindent
  \begin{minipage}[t]{0.52\linewidth}#1\end{minipage}\hfill
  \begin{minipage}[t]{0.46\linewidth}#2\end{minipage}%
  \par\medskip
}
\providecommand{\kienthuccot}[2]{\haicotchay{#1}{#2}}
% Dự phòng nếu chưa nạp graphicx / caption (không ghi đè bản thật)
\providecommand{\resizebox}[3]{#3}
\providecommand{\captionof}[2]{\par\smallskip\noindent\textit{#2}\par}
\newenvironment{hoatdong}[1]{%
  \par\medskip\noindent\textbf{Hoạt động.} #1\par\smallskip
}{\par\medskip}
\newenvironment{emcobiet}{%
  \par\medskip\noindent\textbf{Em có biết}\par\smallskip
}{\par\medskip}
\newenvironment{emdahoc}{%
  \par\medskip\noindent\textbf{Em đã học}\par\smallskip
}{\par\medskip}
\newenvironment{emcothe}{%
  \par\medskip\noindent\textbf{Em có thể}\par\smallskip
}{\par\medskip}
\newenvironment{kienthuc}{%
  \par\medskip\noindent\textbf{Kiến thức cốt lõi}\par\smallskip
}{\par\medskip}
\newenvironment{traloi}{%
  \par\smallskip\noindent\textit{Trả lời / ghi chú của em}\par
  \vspace{1.2em}%
}{\par\medskip}
\makeatother

% từ gốc repo:
%   % Định nghĩa lệnh Lý thuyết — dùng khi biên dịch lt.tex bằng TeX.
% App web đọc lt.tex trực tiếp (bỏ \input file này).
% Trong lt.tex, từ thư mục bài:
%   \input{../../../../_lenh/lythuyet.tex}
% từ gốc repo:
%   \input{ngan-hang/_lenh/lythuyet.tex}
\makeatletter
\providecommand{\indam}[1]{\textbf{#1}}
\providecommand{\chuy}[1]{\par\smallskip\noindent\textit{#1}\par}
\providecommand{\luuy}[1]{\par\smallskip\noindent\textbf{Lưu ý.} #1\par}
\providecommand{\ghichu}[1]{\par\smallskip\noindent\textbf{Ghi chú.} #1\par}
\providecommand{\vidu}[1]{\par\smallskip\noindent\textbf{Ví dụ.} #1\par}
\providecommand{\Blythuyet}{\subsection*{Lý thuyết}}
% Hai cột: kiến thức | hình
\providecommand{\haicotchay}[2]{%
  \par\medskip\noindent
  \begin{minipage}[t]{0.52\linewidth}#1\end{minipage}\hfill
  \begin{minipage}[t]{0.46\linewidth}#2\end{minipage}%
  \par\medskip
}
\providecommand{\kienthuccot}[2]{\haicotchay{#1}{#2}}
% Dự phòng nếu chưa nạp graphicx / caption (không ghi đè bản thật)
\providecommand{\resizebox}[3]{#3}
\providecommand{\captionof}[2]{\par\smallskip\noindent\textit{#2}\par}
\newenvironment{hoatdong}[1]{%
  \par\medskip\noindent\textbf{Hoạt động.} #1\par\smallskip
}{\par\medskip}
\newenvironment{emcobiet}{%
  \par\medskip\noindent\textbf{Em có biết}\par\smallskip
}{\par\medskip}
\newenvironment{emdahoc}{%
  \par\medskip\noindent\textbf{Em đã học}\par\smallskip
}{\par\medskip}
\newenvironment{emcothe}{%
  \par\medskip\noindent\textbf{Em có thể}\par\smallskip
}{\par\medskip}
\newenvironment{kienthuc}{%
  \par\medskip\noindent\textbf{Kiến thức cốt lõi}\par\smallskip
}{\par\medskip}
\newenvironment{traloi}{%
  \par\smallskip\noindent\textit{Trả lời / ghi chú của em}\par
  \vspace{1.2em}%
}{\par\medskip}
\makeatother

\makeatletter
\providecommand{\indam}[1]{\textbf{#1}}
\providecommand{\chuy}[1]{\par\smallskip\noindent\textit{#1}\par}
\providecommand{\luuy}[1]{\par\smallskip\noindent\textbf{Lưu ý.} #1\par}
\providecommand{\ghichu}[1]{\par\smallskip\noindent\textbf{Ghi chú.} #1\par}
\providecommand{\vidu}[1]{\par\smallskip\noindent\textbf{Ví dụ.} #1\par}
\providecommand{\Blythuyet}{\subsection*{Lý thuyết}}
% Hai cột: kiến thức | hình
\providecommand{\haicotchay}[2]{%
  \par\medskip\noindent
  \begin{minipage}[t]{0.52\linewidth}#1\end{minipage}\hfill
  \begin{minipage}[t]{0.46\linewidth}#2\end{minipage}%
  \par\medskip
}
\providecommand{\kienthuccot}[2]{\haicotchay{#1}{#2}}
% Dự phòng nếu chưa nạp graphicx / caption (không ghi đè bản thật)
\providecommand{\resizebox}[3]{#3}
\providecommand{\captionof}[2]{\par\smallskip\noindent\textit{#2}\par}
\newenvironment{hoatdong}[1]{%
  \par\medskip\noindent\textbf{Hoạt động.} #1\par\smallskip
}{\par\medskip}
\newenvironment{emcobiet}{%
  \par\medskip\noindent\textbf{Em có biết}\par\smallskip
}{\par\medskip}
\newenvironment{emdahoc}{%
  \par\medskip\noindent\textbf{Em đã học}\par\smallskip
}{\par\medskip}
\newenvironment{emcothe}{%
  \par\medskip\noindent\textbf{Em có thể}\par\smallskip
}{\par\medskip}
\newenvironment{kienthuc}{%
  \par\medskip\noindent\textbf{Kiến thức cốt lõi}\par\smallskip
}{\par\medskip}
\newenvironment{traloi}{%
  \par\smallskip\noindent\textit{Trả lời / ghi chú của em}\par
  \vspace{1.2em}%
}{\par\medskip}
\makeatother

% từ gốc repo:
%   % Định nghĩa lệnh Lý thuyết — dùng khi biên dịch lt.tex bằng TeX.
% App web đọc lt.tex trực tiếp (bỏ \input file này).
% Trong lt.tex, từ thư mục bài:
%   % Định nghĩa lệnh Lý thuyết — dùng khi biên dịch lt.tex bằng TeX.
% App web đọc lt.tex trực tiếp (bỏ \input file này).
% Trong lt.tex, từ thư mục bài:
%   \input{../../../../_lenh/lythuyet.tex}
% từ gốc repo:
%   \input{ngan-hang/_lenh/lythuyet.tex}
\makeatletter
\providecommand{\indam}[1]{\textbf{#1}}
\providecommand{\chuy}[1]{\par\smallskip\noindent\textit{#1}\par}
\providecommand{\luuy}[1]{\par\smallskip\noindent\textbf{Lưu ý.} #1\par}
\providecommand{\ghichu}[1]{\par\smallskip\noindent\textbf{Ghi chú.} #1\par}
\providecommand{\vidu}[1]{\par\smallskip\noindent\textbf{Ví dụ.} #1\par}
\providecommand{\Blythuyet}{\subsection*{Lý thuyết}}
% Hai cột: kiến thức | hình
\providecommand{\haicotchay}[2]{%
  \par\medskip\noindent
  \begin{minipage}[t]{0.52\linewidth}#1\end{minipage}\hfill
  \begin{minipage}[t]{0.46\linewidth}#2\end{minipage}%
  \par\medskip
}
\providecommand{\kienthuccot}[2]{\haicotchay{#1}{#2}}
% Dự phòng nếu chưa nạp graphicx / caption (không ghi đè bản thật)
\providecommand{\resizebox}[3]{#3}
\providecommand{\captionof}[2]{\par\smallskip\noindent\textit{#2}\par}
\newenvironment{hoatdong}[1]{%
  \par\medskip\noindent\textbf{Hoạt động.} #1\par\smallskip
}{\par\medskip}
\newenvironment{emcobiet}{%
  \par\medskip\noindent\textbf{Em có biết}\par\smallskip
}{\par\medskip}
\newenvironment{emdahoc}{%
  \par\medskip\noindent\textbf{Em đã học}\par\smallskip
}{\par\medskip}
\newenvironment{emcothe}{%
  \par\medskip\noindent\textbf{Em có thể}\par\smallskip
}{\par\medskip}
\newenvironment{kienthuc}{%
  \par\medskip\noindent\textbf{Kiến thức cốt lõi}\par\smallskip
}{\par\medskip}
\newenvironment{traloi}{%
  \par\smallskip\noindent\textit{Trả lời / ghi chú của em}\par
  \vspace{1.2em}%
}{\par\medskip}
\makeatother

% từ gốc repo:
%   % Định nghĩa lệnh Lý thuyết — dùng khi biên dịch lt.tex bằng TeX.
% App web đọc lt.tex trực tiếp (bỏ \input file này).
% Trong lt.tex, từ thư mục bài:
%   \input{../../../../_lenh/lythuyet.tex}
% từ gốc repo:
%   \input{ngan-hang/_lenh/lythuyet.tex}
\makeatletter
\providecommand{\indam}[1]{\textbf{#1}}
\providecommand{\chuy}[1]{\par\smallskip\noindent\textit{#1}\par}
\providecommand{\luuy}[1]{\par\smallskip\noindent\textbf{Lưu ý.} #1\par}
\providecommand{\ghichu}[1]{\par\smallskip\noindent\textbf{Ghi chú.} #1\par}
\providecommand{\vidu}[1]{\par\smallskip\noindent\textbf{Ví dụ.} #1\par}
\providecommand{\Blythuyet}{\subsection*{Lý thuyết}}
% Hai cột: kiến thức | hình
\providecommand{\haicotchay}[2]{%
  \par\medskip\noindent
  \begin{minipage}[t]{0.52\linewidth}#1\end{minipage}\hfill
  \begin{minipage}[t]{0.46\linewidth}#2\end{minipage}%
  \par\medskip
}
\providecommand{\kienthuccot}[2]{\haicotchay{#1}{#2}}
% Dự phòng nếu chưa nạp graphicx / caption (không ghi đè bản thật)
\providecommand{\resizebox}[3]{#3}
\providecommand{\captionof}[2]{\par\smallskip\noindent\textit{#2}\par}
\newenvironment{hoatdong}[1]{%
  \par\medskip\noindent\textbf{Hoạt động.} #1\par\smallskip
}{\par\medskip}
\newenvironment{emcobiet}{%
  \par\medskip\noindent\textbf{Em có biết}\par\smallskip
}{\par\medskip}
\newenvironment{emdahoc}{%
  \par\medskip\noindent\textbf{Em đã học}\par\smallskip
}{\par\medskip}
\newenvironment{emcothe}{%
  \par\medskip\noindent\textbf{Em có thể}\par\smallskip
}{\par\medskip}
\newenvironment{kienthuc}{%
  \par\medskip\noindent\textbf{Kiến thức cốt lõi}\par\smallskip
}{\par\medskip}
\newenvironment{traloi}{%
  \par\smallskip\noindent\textit{Trả lời / ghi chú của em}\par
  \vspace{1.2em}%
}{\par\medskip}
\makeatother

\makeatletter
\providecommand{\indam}[1]{\textbf{#1}}
\providecommand{\chuy}[1]{\par\smallskip\noindent\textit{#1}\par}
\providecommand{\luuy}[1]{\par\smallskip\noindent\textbf{Lưu ý.} #1\par}
\providecommand{\ghichu}[1]{\par\smallskip\noindent\textbf{Ghi chú.} #1\par}
\providecommand{\vidu}[1]{\par\smallskip\noindent\textbf{Ví dụ.} #1\par}
\providecommand{\Blythuyet}{\subsection*{Lý thuyết}}
% Hai cột: kiến thức | hình
\providecommand{\haicotchay}[2]{%
  \par\medskip\noindent
  \begin{minipage}[t]{0.52\linewidth}#1\end{minipage}\hfill
  \begin{minipage}[t]{0.46\linewidth}#2\end{minipage}%
  \par\medskip
}
\providecommand{\kienthuccot}[2]{\haicotchay{#1}{#2}}
% Dự phòng nếu chưa nạp graphicx / caption (không ghi đè bản thật)
\providecommand{\resizebox}[3]{#3}
\providecommand{\captionof}[2]{\par\smallskip\noindent\textit{#2}\par}
\newenvironment{hoatdong}[1]{%
  \par\medskip\noindent\textbf{Hoạt động.} #1\par\smallskip
}{\par\medskip}
\newenvironment{emcobiet}{%
  \par\medskip\noindent\textbf{Em có biết}\par\smallskip
}{\par\medskip}
\newenvironment{emdahoc}{%
  \par\medskip\noindent\textbf{Em đã học}\par\smallskip
}{\par\medskip}
\newenvironment{emcothe}{%
  \par\medskip\noindent\textbf{Em có thể}\par\smallskip
}{\par\medskip}
\newenvironment{kienthuc}{%
  \par\medskip\noindent\textbf{Kiến thức cốt lõi}\par\smallskip
}{\par\medskip}
\newenvironment{traloi}{%
  \par\smallskip\noindent\textit{Trả lời / ghi chú của em}\par
  \vspace{1.2em}%
}{\par\medskip}
\makeatother

\makeatletter
\providecommand{\indam}[1]{\textbf{#1}}
\providecommand{\chuy}[1]{\par\smallskip\noindent\textit{#1}\par}
\providecommand{\luuy}[1]{\par\smallskip\noindent\textbf{Lưu ý.} #1\par}
\providecommand{\ghichu}[1]{\par\smallskip\noindent\textbf{Ghi chú.} #1\par}
\providecommand{\vidu}[1]{\par\smallskip\noindent\textbf{Ví dụ.} #1\par}
\providecommand{\Blythuyet}{\subsection*{Lý thuyết}}
% Hai cột: kiến thức | hình
\providecommand{\haicotchay}[2]{%
  \par\medskip\noindent
  \begin{minipage}[t]{0.52\linewidth}#1\end{minipage}\hfill
  \begin{minipage}[t]{0.46\linewidth}#2\end{minipage}%
  \par\medskip
}
\providecommand{\kienthuccot}[2]{\haicotchay{#1}{#2}}
% Dự phòng nếu chưa nạp graphicx / caption (không ghi đè bản thật)
\providecommand{\resizebox}[3]{#3}
\providecommand{\captionof}[2]{\par\smallskip\noindent\textit{#2}\par}
\newenvironment{hoatdong}[1]{%
  \par\medskip\noindent\textbf{Hoạt động.} #1\par\smallskip
}{\par\medskip}
\newenvironment{emcobiet}{%
  \par\medskip\noindent\textbf{Em có biết}\par\smallskip
}{\par\medskip}
\newenvironment{emdahoc}{%
  \par\medskip\noindent\textbf{Em đã học}\par\smallskip
}{\par\medskip}
\newenvironment{emcothe}{%
  \par\medskip\noindent\textbf{Em có thể}\par\smallskip
}{\par\medskip}
\newenvironment{kienthuc}{%
  \par\medskip\noindent\textbf{Kiến thức cốt lõi}\par\smallskip
}{\par\medskip}
\newenvironment{traloi}{%
  \par\smallskip\noindent\textit{Trả lời / ghi chú của em}\par
  \vspace{1.2em}%
}{\par\medskip}
\makeatother


\subsection{Dạng bài tập mẫu}
\subsubsection{Dạng: Chuyển đổi số đo góc giữa độ và radian}
\begin{phuongphap}
Để chuyển đổi số đo góc giữa độ và radian, ta sử dụng các công thức sau:
\begin{itemize}
    \item Để chuyển từ độ sang radian: $ \alpha_{rad} = \alpha_{độ} \cdot \frac{\pi}{180^\circ} $.
    \item Để chuyển từ radian sang độ: $ \alpha_{độ} = \alpha_{rad} \cdot \frac{180^\circ}{\pi} $.
\end{itemize}
\textbf{Lưu ý:}
\begin{itemize}
    \item Số đo góc theo radian thường không có đơn vị hoặc có đơn vị là "rad".
    \item Số đo góc theo độ có đơn vị là "$^\circ$".
\end{itemize}
\end{phuongphap}
\begin{vidumau}
Đổi số đo của góc $120^\circ$ sang radian. Khẳng định sau đây \textbf{ĐÚNG} hay $\textbf{SAI}$?
\choice
A. $120^\circ = \frac{2\pi}{3}$ rad \True
B. $120^\circ = \frac{\pi}{3}$ rad
C. $120^\circ = \frac{3\pi}{2}$ rad
D. $120^\circ = \frac{2\pi}{5}$ rad
\loigiai{
Để đổi $120^\circ$ sang radian, ta áp dụng công thức: $ \alpha_{rad} = \alpha_{độ} \cdot \frac{\pi}{180^\circ} $.
Thay $ \alpha_{độ} = 120^\circ $ vào công thức, ta được:
$ 120^\circ \cdot \frac{\pi}{180^\circ} = \frac{120\pi}{180} = \frac{2\pi}{3} $ rad.
Vậy, $120^\circ = \frac{2\pi}{3}$ rad.
Khẳng định A là đúng.
}
\end{vidumau}

\subsubsection{Dạng: Xác định góc có cùng điểm biểu diễn trên đường tròn lượng giác}
\begin{phuongphap}
Hai góc lượng giác $\alpha$ và $\beta$ có cùng điểm biểu diễn trên đường tròn lượng giác khi và chỉ khi hiệu của chúng là một bội số nguyên của $2\pi$ (nếu tính theo radian) hoặc $360^\circ$ (nếu tính theo độ).
Tức là:
\begin{itemize}
    \item Nếu tính theo radian: $ \beta = \alpha + k2\pi $, với $k \in \mathbb{Z}$.
    \item Nếu tính theo độ: $ \beta = \alpha + k360^\circ $, với $k \in \mathbb{Z}$.
\end{itemize}
Để kiểm tra hai góc có cùng điểm biểu diễn hay không, ta thực hiện các bước sau:
\begin{enumerate}
    \item Đưa cả hai góc về cùng một đơn vị đo (độ hoặc radian).
    \item Tính hiệu của hai góc đó.
    \item Kiểm tra xem hiệu đó có phải là bội số nguyên của $2\pi$ (hoặc $360^\circ$) hay không.
\end{enumerate}
\end{phuongphap}
\begin{vidumau}
Các cặp góc sau có cùng điểm biểu diễn trên đường tròn lượng giác. Khẳng định sau đây \textbf{ĐÚNG} hay $\textbf{SAI}$?
Cặp góc $ \frac{\pi}{4} $ và $ \frac{9\pi}{4} $ có cùng điểm biểu diễn trên đường tròn lượng giác.
\choice
A. ĐÚNG \True
B. SAI
\loigiai{
Để kiểm tra xem hai góc $ \frac{\pi}{4} $ và $ \frac{9\pi}{4} $ có cùng điểm biểu diễn trên đường tròn lượng giác hay không, ta tính hiệu của chúng:
$ \frac{9\pi}{4} - \frac{\pi}{4} = \frac{8\pi}{4} = 2\pi $.
Vì hiệu của hai góc là $2\pi$, tức là $1 \cdot 2\pi$ (với $k=1 \in \mathbb{Z}$), nên hai góc này có cùng điểm biểu diễn trên đường tròn lượng giác.
Khẳng định A là đúng.
}
\end{vidumau}

\subsubsection{Dạng: Xác định góc phần tư của góc lượng giác}
\begin{phuongphap}
Để xác định góc phần tư của một góc lượng giác $\alpha$, ta thực hiện các bước sau:
\begin{enumerate}
    \item Đưa góc $\alpha$ về dạng $ \alpha_0 + k2\pi $ (hoặc $ \alpha_0 + k360^\circ $), trong đó $0 \le \alpha_0 < 2\pi$ (hoặc $0^\circ \le \alpha_0 < 360^\circ$).
    \item Dựa vào giá trị của $\alpha_0$ để xác định góc phần tư:
    \begin{itemize}
        \item Góc phần tư I: $0 < \alpha_0 < \frac{\pi}{2}$ (hoặc $0^\circ < \alpha_0 < 90^\circ$).
        \item Góc phần tư II: $\frac{\pi}{2} < \alpha_0 < \pi$ (hoặc $90^\circ < \alpha_0 < 180^\circ$).
        \item Góc phần tư III: $\pi < \alpha_0 < \frac{3\pi}{2}$ (hoặc $180^\circ < \alpha_0 < 270^\circ$).
        \item Góc phần tư IV: $\frac{3\pi}{2} < \alpha_0 < 2\pi$ (hoặc $270^\circ < \alpha_0 < 360^\circ$).
    \end{itemize}
    \textbf{Lưu ý:} Các góc $0, \frac{\pi}{2}, \pi, \frac{3\pi}{2}, 2\pi$ (hoặc $0^\circ, 90^\circ, 180^\circ, 270^\circ, 360^\circ$) nằm trên các trục tọa độ và không thuộc bất kỳ góc phần tư nào.
\end{enumerate}
\end{phuongphap}
\begin{vidumau}
Góc $ \frac{7\pi}{3} $ thuộc góc phần tư thứ mấy? Các mệnh đề sau đây \textbf{ĐÚNG} hay $\textbf{SAI}$?
Góc $ \frac{7\pi}{3} $ thuộc góc phần tư thứ I.
\choice
A. ĐÚNG \True
B. SAI
\loigiai{
Để xác định góc phần tư của góc $ \frac{7\pi}{3} $, ta đưa nó về dạng $ \alpha_0 + k2\pi $ với $0 \le \alpha_0 < 2\pi$.
Ta có: $ \frac{7\pi}{3} = \frac{6\pi + \pi}{3} = \frac{6\pi}{3} + \frac{\pi}{3} = 2\pi + \frac{\pi}{3} $.
Ở đây, $ \alpha_0 = \frac{\pi}{3} $ và $k=1$.
Vì $0 < \frac{\pi}{3} < \frac{\pi}{2}$, nên góc $ \frac{7\pi}{3} $ có điểm biểu diễn nằm trong góc phần tư thứ I.
Khẳng định A là đúng.
}
\end{vidumau}

\subsubsection{Dạng: Tính độ dài cung tròn khi biết bán kính và số đo góc}
\begin{phuongphap}
Để tính độ dài cung tròn $l$ khi biết bán kính $R$ và số đo góc ở tâm $\alpha$, ta sử dụng công thức:
$ l = R \cdot \alpha $
Trong đó:
\begin{itemize}
    \item $l$: độ dài cung tròn.
    \item $R$: bán kính của đường tròn.
    \item $\alpha$: số đo góc ở tâm chắn cung đó, \textbf{phải được tính bằng radian}.
\end{itemize}
Nếu số đo góc được cho bằng độ, ta phải chuyển đổi sang radian trước khi áp dụng công thức: $ \alpha_{rad} = \alpha_{độ} \cdot \frac{\pi}{180^\circ} $.
\end{phuphuongphap}
\begin{vidumau}
Cho đường tròn có bán kính bằng $9$ (cm). Tìm số đo (theo radian) của cung có độ dài $3\pi$ (cm).
\choice
A. $ \frac{\pi}{3} $ \True
B. $ \frac{\pi}{2} $
C. $ \pi $
D. $ \frac{2\pi}{3} $
\loigiai{
Gọi $R$ là bán kính đường tròn, $l$ là độ dài cung tròn và $\alpha$ là số đo góc ở tâm (theo radian).
Theo đề bài, ta có:
$R = 9$ cm
$l = 3\pi$ cm
Công thức tính độ dài cung tròn là $l = R \cdot \alpha$.
Thay các giá trị đã biết vào công thức, ta được:
$3\pi = 9 \cdot \alpha$
Để tìm $\alpha$, ta chia cả hai vế cho $9$:
$ \alpha = \frac{3\pi}{9} = \frac{\pi}{3} $ (radian).
Vậy, số đo của cung là $ \frac{\pi}{3} $ radian.
}
\end{vidumau}
